\section{Research Objectives}
\label{sec:research_objectives}

The researchers aim to develop, test, and identify the existence of the Braess Paradox within Iloilo City’s network, specifically the primary roads of the Iloilo City's seven districts: Villa-Arevalo, City Proper, Jaro, La Paz, Lapuz, Mandurriao, Molo. The results will provide a foundation for future city-wide analysis and contribute to more data-driven decisions in road planning and congestion management. Specifically, the study aims to:

\begin{itemize}
    \item Collect and process road network data, including road length, travel time, traffic volume, and other factors in the primary roads of the Iloilo City's seven districts: Villa-Arevalo, City Proper, Jaro, La Paz, Lapuz, Mandurriao, Molo
    \item Develop a transportation network model that represents the current traffic conditions of the study area.
    \item Simulate various road closure and modification scenarios to observe changes in overall travel time and traffic distribution.
    \item Identify and analyze routes affected by the Braess Paradox, in which adding or removing a road segment can counterintuitively improve network efficiency.
    \item Validate the developed model by comparing simulated outcomes with existing literature and real-world observations, ensuring the reliability of the algorithm and results.
    \item Evaluate the effectiveness of road closure strategies in reducing travel time and improving traffic flow in the study area.
\end{itemize}